\section{泛化界的解释力与边界}
\label{sec:13-Machine-Learning/10-Learning-Theory/04}

你把前三节的公式写进设计文档，准备给模型的线上错误率标一个上界。数据十万条，线性模型一百个特征，代入 VC 界，算出来的界宽是十一个百分点。可你的线下测试错误率只有 \(2.3\%\)，界给出的``保证''比实测差了五倍。换成深度网络更难看：参数量一上百万，界宽直接超过 \(1\)，而错误率天生不会超过 \(1\)，这条界什么也没说。

于是必须把话讲全。这些界在数值上几乎从不具备预测力。推导没有写错，它们承担的任务与你此刻想要的任务本来就不同。本节先把松弛的来源摆清楚，再看它们确实解释了什么、解释不了什么，最后交代读者应该带走的判断方式。

\subsection{先承认：这些界在数值上很松}

代一组数字。\(\R^{100}\) 中的线性分类器有 \(d_{\mathrm{VC}}=101\)，取 \(n=10^5\)、\(\delta=0.05\)：

\[
\sqrt{\frac{d_{\mathrm{VC}}\log n+\log(1/\delta)}{n}}
\approx\sqrt{\frac{101\times11.5+3}{10^{5}}}
\approx 0.11.
\]

约十一个百分点，还是忽略常数之后的量级；VC 定理原始形式里的常数会让它再宽好几倍。对深度网络，VC 维至少随参数量增长，代入任何现实的 \(n\)，界宽都会突破 \(1\)。

松弛有三个明确的来源，每一个都是为``保证成立''付的保费。其一，界对所有数据分布一致成立，包括那些专门设计来坑你的分布，而你的数据显然不是最坏的那个。其二，联合界或生长函数按最坏点配置计价，实际数据的几何往往温和得多。其三，界只用到损失有界这一条，没有用到问题的任何结构。把``界大于 \(1\)''读成``理论错了''是误读；这类界的产出是结构信息——泛化间隙随什么收紧、随什么变松——通常不在于给出一个小数值。

\subsection{松归松，它解释了正则化与结构风险最小化}

把假设类组织成嵌套序列 \(\mathcal{H}_1\subset\mathcal{H}_2\subset\cdots\)，记 \(d_k=d_{\mathrm{VC}}(\mathcal{H}_k)\) 随 \(k\) 递增，对每个类各写一条界：

\[
R(h)\le\widehat R(h)+\mathrm{width}\bigl(d_k,n,\delta\bigr),
\qquad h\in\mathcal{H}_k.
\]

大类里 \(\widehat R\) 能压得更低，界宽却更大；小类反过来。在 \(\mathcal{H}_k\) 上做 ERM 得到 \(\widehat h\)，把它的超出风险按类内最优 \(h_k^*\in\argmin_{h\in\mathcal{H}_k}R(h)\) 与全问题最优风险 \(R^*\) 拆开：

\[
R\bigl(\widehat h\bigr)-R^*
=\underbrace{R\bigl(\widehat h\bigr)-R(h_k^*)}_{\text{估计误差}}
+\underbrace{R(h_k^*)-R^*}_{\text{逼近误差}}.
\]

类越大，逼近误差越小，估计误差的界越宽。结构风险最小化（SRM）把这个权衡摆到明面上：在嵌套序列上最小化``经验风险加界宽''，让数据自己挑复杂度档位，不必预先固定一个类。惩罚型目标

\[
\min_{h}\;\widehat R(h)+\lambda\,\Omega(h)
\]

是它的连续版本。\(\Omega\) 大意味着有效类大，\(\lambda\) 就是为泛化间隙标出的价格。\hyperref[sec:09-Convex-Optimization/07-Applications/01]{``光滑经验风险与 \(L_2\) 正则化''}中正则项的统计角色到这里才补齐：它一边改善条件数，一边为复杂度付费，偏差与界宽的权衡被写进了目标函数本身。\hyperref[sec:11-Mathematical-Statistics/06-Resampling-and-Model-Selection/04]{``偏差---方差、正则化与模型选择''}从估计量的方差一侧给出同一权衡的另一种记法，两边可以对照着读。

\subsection{它也解释了 SVM 为什么盯间隔不盯维度}

间隔界换了一个复杂度来源。若数据落在半径 \(B\) 的球内，且分类器以几何间隔 \(\gamma\) 把两类分开，则有效复杂度的量级是 \((B/\gamma)^2\)，与输入维数无关——即使核方法把数据映到无穷维特征空间，界照常成立。\hyperref[sec:09-Convex-Optimization/07-Applications/03]{``最大间隔与支持向量机''}把目标写成``固定函数间隔、最小化 \(\norm{w}\)''，统计理由正在这里：在那个约定下几何间隔等于 \(1/\norm{w}\)，所以 SVM 的目标函数本身就是一条泛化界的两项，hinge 项管 \(\widehat R\)，\(\norm{w}^2\) 项管界宽。

沿这条路走得更远的是范数界与 Rademacher 复杂度：用``类在这批数据上拟合随机噪声的平均能力''代替最坏情形的组合计数，常常紧得多，而且可以直接从数据估计。它们的推导已经能撑起一门课，这里只登记位置——当 VC 界太松而间隔不适用时，这是下一站。

\subsection{三笔账要分开记：逼近、估计、优化}

上面的分解只有两项，实际的训练流程还漏了一项。你手上的 \(\widehat h\) 未必是 \(\argmin_{h\in\mathcal{H}}\widehat R(h)\)，它只是优化算法在有限步之后停下的地方。把三笔账写全：

\[
R\bigl(h_{\text{输出}}\bigr)-R^*
=\underbrace{R(h_k^*)-R^*}_{\text{逼近}}
+\underbrace{R\bigl(\widehat h\bigr)-R(h_k^*)}_{\text{估计}}
+\underbrace{R\bigl(h_{\text{输出}}\bigr)-R\bigl(\widehat h\bigr)}_{\text{优化}}.
\]

\hyperref[sec:09-Convex-Optimization/08-Nonconvex-and-Stochastic-Optimization/06]{``优化误差、统计误差与泛化不能混为一谈''}把这三项的相互制约讲得更细：优化误差降到远小于统计误差之后继续降就没有收益，而在非凸地形上把训练损失优化到极致，往往还会顶高第二项。诊断一个泛化不佳的模型，第一步是判断落差主要挂在哪一项上——训练误差本身就高，问题多半在逼近或优化；训练误差近零而测试误差高，问题在估计一侧。三项混作一谈时，调参就会退化成碰运气。

\subsection{深度网络的泛化落在经典框架之外}

也必须直面框架失效的地方。以下三条是已被反复复现的实验事实。过参数化网络的参数量远超样本量，按参数计数得到的 VC 维极其巨大；同一批网络能把纯随机标签拟合到零训练误差，也就是完全记住噪声；而同一架构、同一优化器在真实数据上泛化良好。按经典一致收敛的说法，能记住随机标签的类应当毫无泛化保证，这三条事实放在一起，经典框架给不出解释。

双下降进一步撕开了经典图景。测试误差随模型规模先降后升，在插值阈值——容量刚好大到能把训练误差压到零的那一点——附近冲出一个尖峰；越过阈值继续增大模型，误差再次下降，有时降到比第一个谷底还低。

\begin{center}
\begin{tikzpicture}[>=stealth]
\draw[->] (0,0) -- (7.8,0) node[below right] {\footnotesize 模型容量};
\draw[->] (0,0) -- (0,4.8) node[above left] {\footnotesize 误差};
% classical bound width
\draw[dashed,gray] plot[smooth] coordinates {(0.4,0.5) (2.0,1.3) (3.6,2.3) (5.2,3.4) (6.9,4.5)};
\node[gray,anchor=south east] at (6.9,4.45) {\footnotesize 经典界宽};
% test error: double descent
\draw[very thick] plot[smooth] coordinates {(0.4,3.6) (1.2,2.6) (2.0,2.05) (2.6,1.95) (3.1,2.5) (3.5,3.9) (4.1,2.6) (5.0,1.85) (6.0,1.5) (6.9,1.35)};
\node[anchor=west] at (6.95,1.35) {\footnotesize 测试};
% train error
\draw[thick,gray!75] plot[smooth] coordinates {(0.4,3.2) (1.2,2.2) (2.0,1.4) (2.8,0.7) (3.5,0.15) (4.5,0.1) (6.9,0.08)};
\node[anchor=west,gray!75] at (6.95,0.35) {\footnotesize 训练};
% interpolation threshold
\draw[dotted] (3.5,0) -- (3.5,3.9);
\node[anchor=north] at (3.5,-0.05) {\footnotesize 插值阈值};
\node[anchor=south] at (1.9,3.9) {\footnotesize 经典 U 形};
\node[anchor=south] at (5.6,2.5) {\footnotesize 经典框架未覆盖};
\end{tikzpicture}
\end{center}

\emph{插值阈值左侧的 U 形与经典权衡吻合，右侧测试误差随容量继续下降，而经典界宽在整条横轴上单调上升。}

图里那条单调上升的虚线是关键。经典界只知道容量变大、保证变松，它没有任何机制描述阈值右侧发生的事。目前较受支持的线索是：有效复杂度不只由假设类决定，还由优化过程决定。初始化、SGD 的隐式偏置、早停、数据本身的低维结构，都在把实际搜索范围收进 \(\mathcal{H}\) 的某个``好''子集，而这个子集的复杂度远小于整个类。隐式正则化、神经正切核、基于算法稳定性的界都在改写``复杂度从哪里来''这个问题。

需要标清楚的是，上一段的解释目前仍属猜想。哪些机制起主导、能否给出对真实网络数值有效的界、双下降的尖峰在多大程度上由标签噪声驱动，都还没有公认答案。已知的是那三条实验事实和经典界在此失效；未知的是替代框架长什么样。把这一段当作前沿记住，而不是当作教材里的缺陷。

\subsection{带走的判断方式}

泛化要求有效复杂度相对样本量受控，而``模型简单''这句话太粗，没法当作判断依据。控制的旋钮装在好几个不同的地方：样本量让所有界按 \(\sqrt{n}\) 收紧；假设类的大小以 \(\log\abs{\mathcal{H}}\) 或 \(d_{\mathrm{VC}}\) 计价；几何量以间隔 \((B/\gamma)^2\) 或范数 \(\norm{w}\) 计价；优化过程的隐式偏置目前落在经典框架之外，却实实在在地在现代实践里起作用。

这条理论链也把早已见过的工程纪律串成了推论：测试集只能用一次，调参必须走验证集或重抽样（\hyperref[sec:11-Mathematical-Statistics/06-Resampling-and-Model-Selection/03]{``交叉验证评估的是完整训练流程''}），训练损失下降从来不是会泛化的证据。理论交给你的是一张核对清单，而非一个可以填进设计文档的数字：每当数据参与一个决定，就问一次——这次决定花掉了多少随机性预算，它对应界里的哪一项。清单之外的部分，仍然要靠留出的那批数据说话。
